%\vspace{1.5\baselineskip}
\begin{figure}[H]
    \begin{center}
        \begin{circuitikz}[european]
            \ctikzset{tripoles/mos style/arrows}[line width=thick, european]
            
            % Transistoren 
            \draw (-2,0) node[nmos, yscale=-1, rotate=180](nmosL){}
                node at (nmosL.center) [anchor=center, xshift=-10pt]{$N_0$};
            \draw (2,0) node[nmos](nmosR){$N_1$};

            \draw (nmosL.gate) -- (nmosR.gate)
                coordinate[pos=0.5](gate_cross)
                (gate_cross) node[circle, fill, inner sep=1.5pt]{};

            \draw (nmosL.drain) -- ++(0,0.25) 
                coordinate(left_cross)
                node[circle, fill, inner sep=1.5pt]{};

            \draw (left_cross) -| (gate_cross);

            \draw (left_cross) -- ++(0,1.25);
            \draw (nmosR.drain) -- ++(0,1.5);

            % Ground 
            \draw (nmosL.source) -- ++ (0,-0.5) -- ++ (2,0)
                node[circle, fill, inner sep=1.5pt]{};
            \draw (nmosR.source) -- ++ (0,-0.5) -- ++ (-2,0)
                coordinate(ground);
            \draw (ground) -- ++(0,-.25) node[sground]{}
                node[right, xshift=5pt, yshift=-7pt]{Gnd};

            % Ströme
            \draw[->, thick, >=stealth] (-2.25,2) -- ++(0,-.75)
                node[left, yshift=10pt]{$I_{in}$};
            \draw[->, thick, >=stealth] (2.25,2) -- ++(0,-.75)
                node[right, yshift=10pt]{$I_{out}$};

        \end{circuitikz}
    \end{center}
    \caption[Aufbau Stromspiegel]{\\Quelle: Eigene Darstellung Aufbau Stromspiegel in Anlehnung an \cite[S. 512]{MicroelectronicCircuits}}
    \label{fig:current_mirror}
\end{figure}